\clearpage
\section{Derivation and Properties of the \texorpdfstring{$\kappa$}{kappa} Extension}
\label{app:kappa-derivation}
\setcounter{equation}{0}\setcounter{table}{0}\setcounter{figure}{0}
The extension changes assignment while retaining the baseline primitives in Section~\ref{sec:monte-carlo-design}. The excluded sum has variance two, $\epsilon_i$ has variance one, and $\operatorname{Cov}(u_i,\epsilon_i)=0.3$. Introduce independent $w_i\sim\mathcal N(0,1)$ and write
\begin{equation}
 d_i^*(\kappa)=\kappa(z_{1i}+z_{2i})+c(\kappa)w_i+\epsilon_i,
 \qquad \kappa\in[0,1].
 \label{eq:appendix-kappa-unknown}
\end{equation}
Independence eliminates all cross-covariances between index components. Preserving the baseline variance requires
\[
2\kappa^2+c(\kappa)^2+1=3,
\qquad c(\kappa)=\sqrt{2(1-\kappa^2)},\quad 0\le\kappa\le1.
\]
The nonnegative root is conventional; normal symmetry makes its sign distributionally immaterial. This coefficient is specific to the independent-noise construction. At $\kappa=0.50$, the signal variance is 0.5, replacement variance is 1.5, and $c(0.50)=\sqrt{1.5}\approx1.2247$. Calculations use the unrounded coefficient.
\begin{equation}
 \begin{aligned}
 d_i^*(\kappa)&=\kappa(z_{1i}+z_{2i})+\sqrt{2(1-\kappa^2)}\,w_i+\epsilon_i,\\
 D_i(\kappa)&=\Phi\!\left(d_i^*(\kappa)\right).
 \end{aligned}
 \label{eq:appendix-kappa-assignment}
\end{equation}
Normality gives $d_i^*(\kappa)\sim\mathcal N(0,3)$, where 3 denotes variance. Thus for $0<d<1$,
\begin{equation}
 \begin{aligned}
 \Pr\!\left[D_i(\kappa)\le d\right]
 &=\Pr\!\left[\Phi(d_i^*(\kappa))\le d\right]\\
 &=\Pr\!\left[d_i^*(\kappa)\le\Phi^{-1}(d)\right]\\
 &=\Phi\!\left(\frac{\Phi^{-1}(d)}{\sqrt{3}}\right).
 \end{aligned}
 \label{eq:appendix-treatment-cdf}
\end{equation}
The CDF is zero below zero and one above one. It is independent of $\kappa$, but is not Uniform$(0,1)$. Equal variance establishes this invariance because the index is normal, not for arbitrary distributions.

Write $Z_{ji}=z_{ji}+v_{ji}+C_{ji}$, where $C_{1i}=X_{2i}+X_{3i}+X_{4i}$ and $C_{2i}=X_{7i}+X_{8i}+X_{9i}+X_{10i}$. Independence yields
\[
\operatorname{Cov}(Z_{ji},d_i^*(\kappa))
=\kappa\operatorname{Cov}(z_{ji},z_{1i}+z_{2i})=\kappa,
\qquad j=1,2.
\]
The nonzero means of the control sums do not alter this covariance calculation. It concerns the latent index, and does not assert $\operatorname{Cov}(Z_{ji},D_i(\kappa))=\kappa$. Designed relevance is therefore distinct from a formal IVQR identification statistic.

Similarly, $\operatorname{Cov}(u_i,d_i^*(\kappa))=\operatorname{Cov}(u_i,\epsilon_i)=0.3$. Joint normality gives
\begin{equation}
 \begin{pmatrix}u_i\\d_i^*(\kappa)\end{pmatrix}
 \sim\mathcal{N}\!\left(
 \begin{pmatrix}0\\0\end{pmatrix},
 \begin{pmatrix}1&0.3\\0.3&3\end{pmatrix}\right).
 \label{eq:appendix-endogeneity-invariance}
\end{equation}
The invariant mean vector and covariance matrix establish invariance of the entire joint distribution, hence of $(u_i,\Phi(d_i^*(\kappa)))$. Covariance alone would not suffice without joint normality. Reusing primitive draws aids matched comparisons but is not needed for this distributional result; individual treatment realizations can differ across relevance settings.

At fixed $d>0$, $Y_i(d)=1+5\sum_{j=1}^7X_{ji}+d(1+u_i)$. Independence of $u_i$ and $X_i$ gives $Q_{Y_i(d)\mid X_i}(\tau)=1+5\sum_{j=1}^7X_{ji}+d[1+\Phi^{-1}(\tau)]$. The unchanged potential-outcome equation therefore preserves $\alpha_0(\tau)$. This is a structural quantile at fixed treatment, not an observational conditional quantile at realized endogenous treatment.

At $\kappa=1$ the replacement noise vanishes and the baseline is recovered. At the conceptual endpoint $\kappa=0$, $d_i^*(0)=\sqrt2w_i+\epsilon_i$, so the excluded primitive signal disappears. The weakest simulated setting retains signal variance $2(0.10)^2=0.02$. These endpoint statements establish properties of this construction, not global identification results for arbitrary IVQR models.

Table~\ref{tab:kappa-illustration} records the analytical variance allocation for all simulated settings.
\begin{table}[!htbp]
\caption{Variance allocation in the $\kappa$-extension.}
\label{tab:kappa-illustration}
\centering
\begin{tabular}{rrrrr}
\toprule
$\kappa$ & \shortstack{Coefficient\\$\sqrt{2(1-\kappa^2)}$} & \shortstack{Signal variance\\$2\kappa^2$} & \shortstack{Noise variance\\$2(1-\kappa^2)$} & \shortstack{Total index\\variance} \\
\midrule
1.00 & 0.0000 & 2.0000 & 0.0000 & 3.0000 \\
0.50 & 1.2247 & 0.5000 & 1.5000 & 3.0000 \\
0.25 & 1.3693 & 0.1250 & 1.8750 & 3.0000 \\
0.10 & 1.4071 & 0.0200 & 1.9800 & 3.0000 \\
\bottomrule
\end{tabular}
\begin{tablenotes}
Values are calculated directly from the construction, not simulated. The disturbance $\epsilon_i$ contributes the remaining variance of one in every row.
\end{tablenotes}
\end{table}

\FloatBarrier
